\begin{abstract}
\large
Urban Scout is a mobile-friendly web application designed to reduce the fragmentation of travel information for tourists navigating unfamiliar cities. In many real-world situations, travelers must switch between multiple platforms to access transit schedules, restaurant information, landmarks, local events, and safety updates. This creates inefficiency, confusion, and unnecessary cognitive overload, especially for short-term visitors who need quick and reliable access to location-based information. Urban Scout addresses this problem by integrating these services into a single relational database-backed system.

\vspace{1em}
The project was implemented using a MySQL database, an Express/Node.js backend, and a React frontend with a mobile-first interface. The system supports multiple user roles, primarily tourists and administrators. Tourists can browse places by region, view transit-related information, read localized news and safety alerts, submit incident reports, and manage itineraries. Administrators can manage news content and verify reported incidents. The database design emphasizes relational connections between regions, places, transit services, news posts, alerts, users, and itineraries, allowing the application to present connected information that would normally be spread across different platforms.

\vspace{1em}
Overall, Urban Scout demonstrates how a well-designed relational database can support a unified user experience for tourism-related planning and awareness. By combining transportation, destination browsing, local updates, and trip planning into one system, the project reduces app fatigue and improves usability for travelers on the go.
\end{abstract}
